\section*{Acknowledgements}
\addcontentsline{toc}{section}{Acknowledgements}

Plan of everyone I need to thank, to be written into prose:

Dr. Aaron Knoll, for being my supervisor
- Grounding me on my progress and helping me along all the way through admin and safety
- Repeatedly told me where I actually was, like telling me to get an extension

Dr. Lubbock
- Lending me the torquemeter

Mark Grant
- Guiding me on mechanical design, allowing me to test in the hypersonic lab


Hackspace:
Fred Holdsworth, David Miller, Mohammed Mehdi Nasr Esfhani
- teaching me all about machining over the past 4 years, hundreds of hours of guidance and especially so in this project

Nigel
- Guiding me on safety, working through and being patient with this project with a lot of risk and working through the incident report

David Wright
- Lending me a lot of fittings from the thermofluids lab of mecheng

Franco
- Lending me a lot of the electronics equipment, the box, safety relay, a lot of guidance on safety wiring

Xiao
- Giving advice and teaching me about the magtrol system

Liam Donelly
- Rocket turbine FYP who I borrowed a lot of knowledge and equipment from, guiding me on magtrol and the general direction of the project

Rui, Safee
- Staying by my side as I test and occasionally even helping me set up and debug. Rui especially helped me solve many many problems

Aleks Kent
- Starting the E-pump project in 2023 and inspiring me to get into turbomachinery, and being my best friend

Jordan Farrar, Mark Thornton, Ian Pardew, 
- For helping me 3d print and laser cut parts, general manufacturing

Roland Hutchins
- Allowing me to borrow the garden pump

ICLR - Opportunity and wehere I learned everything
Martin England, Elliot Brookes Gayton, Kiran de Silva
 - For bringing me into this world of rocketry, electronics and propulsion when I was a first year

Airborne engineering
- where I learned a lot a lot of test rig and practical skills

Rudra Nakade
- My great friend who I learned so much about doing rocketry with while working with him, and being a poweruser, who eventually convinced me to do a turbopump FYP out of my many rocketry interests

Elspeth Lau
- My great friend who supported me during the harder times of this project

Mom and Dad
- For supporting me all I life, sacrificing so much of their life to get me into imperial and supporting my life dream of being a rocket engineer, always being there for me and cheering me on

Mario
- Guidance on pump design, and reading this posts on r/rocketry discord let me learn a lot

Oliver Esmond, Tim Barrett
- allowing me to borrow ICSS equipment

Chris Noon
- allowing me to use their swagelok valvse and test rig equipment from turbomachinery lab

Tian
- teaching me about mechanical from 2022 and being my great friend ever since


\section{Introduction}

Rocket engines require high pressures to operate efficiently. Small scale liquid propulsion such as satellite propulsion and amateur liquid roket propulsion runs almost entirely on pressure fed feed systems due to their simplicity and reliability. Larger orbital rockets use turbopumps because the mass of the required tanks becomes prohibative. However, as space launch becomes exponentially more accessible, there is a growing interest in small scale rocket propulsion systems. For example, lunar landers, satellite propulsion and amateur liquid rockets are seeing a resurgence in interest. Due to the historical focus on large turbopumps for such applications, turbopumps at the smaller scale are not as well documented and have different design considerations. A turbopump is a pump and the turbine that drives it, and both are punished by the small scale.

\subsection{Small scale rocket pump challenges}

Pump driven propulsion is the solution orbital liquid rockets use to solve the tank mass problem. By storing propellents at lower pressures and pumping them to the required pressures allows for a large reduction in tank mass. The pump is usually driven by turbines powered by burning a portion of the propellents, or by the heat generated by the cooling of the engine. In recent years, battery and electric pump technology has advanced to such a point that for smaller rockets such as the Electron rocket by Rocket Lab \cite{electron}, electric pumps are starting to become a more viable alternative to turbine driven pumps. \todo{this paragraph might be redundant?}

Pumps are generally characterized by their flow rate $Q_{\textrm{opt}}$, the head $H_{\textrm{opt}}$ and rotor speed $n$ \cite{gulich}. These three parameters largely decide the pump geometry and type of pump. These three parameters can be related by a quantity known as specific speed, $n_q$, $N_s$ or $\omega_s$ depending on the units and definitions used:

\begin{equation}
    n_q = \frac{n \sqrt{Q}}{H^{3/4}} \quad [\text{RPM}, \text{m}^3/\text{s}, \text{m}] \label{eq:specific_speed}
\end{equation}

Different pump geometries are optimal for different regions of specific speed as can be shown by \cref{fig:pump_type_by_specific_speed}. Most liquid rocket engines operate between a specifc speed $n_q$ of 20 to 40 and utilize the centrifugal radial impeller to maximize efficiency. This works great for large rocket engines, but scaling this down to smaller engines becomes increasingly difficult. All rocket engines require high pressures, but smaller rocket engines require less flow rate. To stay in the specific speed according to \cref{eq:specific_speed} that allows for high efficiency in the centrifugal pump requires increasing rotational speed.

To achieve the specific speed of \valSamplepumpNq for a smaller rocket engine of around \valSamplepumpThrust requiring \valSamplepumpMdot of kerosene at \valSamplepumpP leads to rotational speeds of around \valSamplepumpN . Such high rotational speeds leads to challenging mechanical design conditions. To solve this problem, a single pump can be separated into multiple stages, separating the head rise to increase the specific speed of each stage. This is a common solution to pumping hydrogen. This thesis however addresses the other solution, which is to deal with the lower specific speed.

As shown in \cref{fig:pump_type_by_specific_speed}, at lower specific speeds, the partial emission pump becomes the optimum efficiency pump geometry at specific speeds of \todo{get a real number here} 5-10. However, even though these pumps are the most efficient at these low specific speeds, they physically cannot achieve the high efficiencies of a full emission pump. This is the main limitation of trying to operate at low specific speeds.  \todo{maybe not mention full/partial emission so early?}

Due to its lower efficiency and the lack of pump fed liquid rockets for small engines, there is a lack of literature, design knowledge and experimental data surrounding these pumps. This thesis aims to address this gap by reviewing existing literature, designing and manufacturing a partial emission centrifugal pump and comparing its experimental performance specifically in the use case of rocket engines.

\subsection{Small scale rocket turbine challenges}

On the other end of the turbopump shaft is the turbine which drives the pump. The turbine faces the same problem as the pump. If using the gas generator cycle, the small turbine leads to lower blade speeds $U$. To extract the maximum work from from the gas, gas velocity $C_0$ must be high, leading to a supersonic blade turbine. The high gas velocity relative to blade speed leads to the turbine running at a lower efficiency than the optimum $U/C_0 \approx 0.5$ \cite{findsource}. The smaller mass flow also cannot fill the annulus at a sensible blade height, leading to partial admission where only a fraction of the annulus has gas flowing through it at any given time. This makes the stage lossy and awkward to design. In addition to these design limitations, supersonic turbine design is far less studied than subsonic, where methods are scattered across older sources. \todo{need to address subsonic turbines?}

This thesis tackles the design of a turbopump in the context of small scale liquid rocket engines which forces both pump and turbine off optimum design conditions, then demonstrates their coupled performance in a full integration test. The thesis is framed around the design of a \valPumpRpmDesign turbopump using water as a inert simulant for the pumped rocket propellents. Then, a supersonic partial admission impulse turbine is designed to drive the pump using compressed air as the working fluid. The two are then coupled together and tested as a single system to evaluate performance.

\subsection{Aims and Objectives}
The aim of this thesis is to investigate the performance of a small scale rocket low specific speed turbopump with a Barske pump and a supersonic partial admission turbine. To this end, the objectives of this thesis are:

\begin{enumerate}
    \item Pump characterisation: Determine the head-flow curve, efficiency, and cavitation behaviour of a Barske-type impeller across operating speeds
    \item Turbine design and model validation: Design a single-stage partial-admission supersonic impulse turbine with a first-principles computational toolchain, and validate those models against the achievable test data.
    \item Turbopump integration: Demonstrate coupled turbopump operation at scaled off-design conditions and perform analytical extrapolation to the full \valPumpRpmDesign gas generator regime.
    \item Computational toolchain assessment: Assess the 1D/analytical design toolchain against the experimental data, and establish where analytical prediction is sufficient for this class of machine.
\end{enumerate}

\subsection{Scope and Limitations}

This thesis covers the theory, design, manufacture and testing of a small scale turbopump system on inert propellents. The pump was tested up to \valPumpRpmMaxreached, significantly lower than the design speed of \valPumpRpmDesign in the first and final iterations of the conducted experiments. This means the data validation can only happen for the off-design point and extrapolation to the design point is required. Cryogenic and hot gas design and testing are strictly outside the scope of tis thesis. Operational safety procedures are included in the appendix.

\subsection{Thesis Structure}

Chapter 2 establishes the pump theoretical background: pump specific speed, Barske impeller characteristics, impulse turbine velocity triangle theory, and the relevant prior literature. Chapter 3 details the pump design and sizing methodology. Chapter 4 covers turbine design, including structural analysis of the rotor at \valPumpRpmDesign. Chapter 5 describes the custom instrumentation and data acquisition system. Chapter 6 documents the test rig and assembly. Chapter 7 presents the experimental results. Chapter 8 discusses the findings in the context of design predictions and prior literature. Chapter 9 draws conclusions and recommends future work.

\section{Literature and Theory}

\todo{"recent studies" section for pumps and turbines, in particular, MOC is incredibly old and sudhof, erikson have modernized, state this.}

\subsection{Pump Literature and Theory}

\todo{I just realized I don't have the numbers to back me up that disc friction is better. I will waffle, then fill with literature and numbers later on.}

Centrifugal pumpss have multiple main sources of losses. The one that becomes significant at low specific speeds is the disc friction losses. When a rotor object rotates in a fluid, the shear stresses in the fluid cause a drag torque on the object. In most centrifugal pumps the impeller is disc shaped, and the loss is called disc friction loss. \todo{gulich pdf pg 144 says axial case cleareance only has a very small effect on disc friction. It also says that due to leakage flow, partload recirculation and exchange of momentum, tolerance of calculated disc friction loss can be at least +/-25\%, also, somehow head coefficient is related to disc friction? it says to maximize it for low specific speed}. With low specific speed pumps, disk frction is the main source of loss. At $n_q=10$, $\psi_{\textrm{opt}=1}$ and $\mathrm{Re}=10^7$, the disc friction is approxiomately $30\%$ of the usefulf power of the impeller \cite{gulichpage103}. Gulich suggests that a high head coefficient \todo{high head coeffient not defined yet} reduces the diameter and consequently less disk frciont since $P_{RR}\sim D^5$. This is the reason why low specific speed centrifugal pumps physically cannot achieve the same high efficiencies as higher specific speed pumps. \todo{this paragraph is almost fully copied from gulich, need to rephrase.}

\todo{gulich page 429 says impeller losses are low at low specific speeds, referencing fig 3.33. Oop, fig3.33 shows the disc friction loss i think?}On the other hand, at low specific speeds, impeller losses are low as shown by \cref{fig3.33gulich}. \todo{explain this graph} As a result, surface finish on the impeller blades are largely unimportant, allowing for simpler manufacturing processes. Conversely, collector losses become dominant. \todo{need to explain what a collector is} Gulich explains \cite{gulich7.12} that for $n_q<12$ even higher efficiencies are expected than with volutes. This simplifies the design and manufacturing of the pump, which is a significant advantage at small scales where tolerances are both harder to achieve but contribute to more significant losses. \todo{need to explain what a volute is} This is important as the high fluid velocity at the peripheri makes the surface finish of the annular casing of primary importance for pump efficiency \cite{lockdesignfeatures}.

The Barske impeller pump design which was developed for the need for manufacturing simplicity and low specific speed applications became one of the highest efficiency designs at this specific speed. \todo{i might want a section on explaining basic pump and turbine geometry and terminology before this.} The barske impeller operates differently to a normal full emission pump. In a conventional centrifugal pump, impeller blade geometry aims to perfectly guide fluid flow along the blade surfaces in a smooth spiral streamline out of the impeller into the volute \cite{something}. The barske impeller is instead designed to induced almost a pure forced vortex flow where $u=\omega r$ and the fluid is essentially rotating as a solid body with the impeller \cite{barskeunconventionalpaper}. This also explains the unimportance of surface finish on the blades, as low relative velocity between the fluid and the impeller means lower friction losses. 

\todo{I know that conventional pumps are impossible to design at this specific speed, but still need to explain lock's assertion that "The efficiencies of forced vortex pumps in the same specific speed range and capacity are higher than those of similar conventional pump"}

Barske comments that the number of impeller blades does not seem to be critical, with 3-6 being operated successfully. However, altering the number of blades can be used to elminite fluid vibrations \cite{barske1953}. Lock provides corrections for blades from 1 to 16 \cite{lockfig10A}.

The open impeller choice is to further reduce disc friction losses \todo{is this true?} by minimizing surface area. Conventional impellers uses shrouds to guide the flow and reduce secondary flows. However, as in a Barske pump the impeller serves only to generate the forced vortex flow, this was deemed unnecessary by at least Barske and Lock. This also removes axial thrust loads on the impeller, helping with sealing and mechanical design of the shaft. Also, simplify manufacturing.

Conventional pump design with current design tools produce nonsensical geometries and results at low specific speeds because of one number that is used in the design process, which is the slip coefficient. The slip coefficient determines the angular deflection of the fluid leaving the impeller from the blade angle at the impeller exit. This number is usually determined empirically based on regression fits from existing data. However, at $n_q\approx 7$, the equations relating blade exit angle, slip coefficient and head coefficient mathematically does not have a converging root. Physically, there is no blade geometry that can bend the blade enough to keep the flow attached to the blade. This is further motivates the forced vortex Barske architecture, as it does not rely on attached flow, removing slip factor from the design process entirely.

All theories relating to pump and turbine design can be derived from the conservation of angular momentum by Newton's second law of motion. In turbomachinery, this is usually expressed in terms of the Euler turbine equation:

\begin{equation}
    P=T_{shaft}\omega=\rho Q \left(r_2 c_{2u} - r_1 c_{1u}\right)\omega=\rho Q\left(u_2c_{2u}-u_1c_{1u}\right)
    \label{eq:euler_turbine_equation}
\end{equation}

Rewriting this equation in terms of specific work $Y=P/\rho Q$ per unit mass:

\begin{equation}
    Y = \left(r_2 c_{2u} - r_1 c_{1u}\right)\omega = u_2c_{2u}-u_1c_{1u}
     \label{eq:euler_specific_work}
\end{equation}

derive from bernouli, with rotating frame substitution $c^2=w^2-u^2+2\times u \times c_u$ and neglecting compressibility losses:

$$p_{1}+\frac{\rho}{2}w_{1}^2-\frac{\rho}{2}u_{1}^2=p_{2}+\frac{\rho}{2}w_{2}^2-\frac{\rho}{2}u_{2}^2$$

and as such the theoretical rise in static pressure:

\begin{equation}
    \Delta p_{st} = p_2-p_1 = \frac{\rho}{2}\left(w_1^2-w_2^2\right)+\frac{\rho}{2}\left(u_2^2-u_1^2\right)
    \label{eq:pressure_rise}
\end{equation}

In Barske, Lobanoff and Lock the relative velocity $w$ is assumed to be negligable, ending up with:

\begin{equation}
    \Delta p_{st} = \frac{\rho}{2}\left(u_2^2-u_1^2\right)
    \label{eq:pressure_rise_negligible_w}
\end{equation}

Then, the kinetic energy increase assuming negligible inlet prerotation is:

\begin{equation}
    \Delta p_{st} = \frac{\rho}{2}u_2^2
    \label{eq:pressure_rise_negligible_w_prerotation}
\end{equation}




\cref{eq:pressure_rise_negligible_w_prerotation} is the theoretical total pressure a barske impeller generates at the blade tips. In a barske pump, the diffuser is responsible for converting this kinetic energy into static pressure by slowing down the fluid in an expanding diffuser. Barske introduces a loss coefficient for the kinetic energy which can be recovered through the diffuser to be between 0.2 and 0.5. Barske uses the symbol $\psi$ for this coefficient, but as it conflicts with the commonly used symbol for head coefficient, this paper will instead use $c_k$. Combining the static and kinetic pressure rise, the total increase in pressure is:

\begin{equation}
    p_t = p_{s_{th}} + c_k p_{d_{th}}=\frac{1}{2}\rho\left((1+c_k) u_2^2 - u_1^2\right)
    \label{eq:barske_pressure_rise}
\end{equation}

This forms a simple basic model for barske pump performance. 

\subsection{Barske Pump sizing}

The barske pump is sized by the methods provided by barske \cite{barske1953} with the assitance of Gulich. First the inlet is sized by choosing a reasonable inlet velocity and calculating the inlet area:

\begin{equation}
    d_0 = \sqrt{\frac{4\dot{m}}{\pi \rho c_0}}
    \label{eq:barske_inlet_diameter}
\end{equation}

The impeller blade inlet diameter is enlarged by a small amount as some pre-rotation is desirable for increased head. For example, $d_1\sim1.1d_0$.

By substituting $u=\pi n d / 60$ into \cref{eq:barske_pressure_rise}, the required impeller tip diameter can be calculated:

\begin{equation}
    d_2 = \sqrt{\frac{1}{1+c_k}\left(\frac{7200p_t}{\rho n^2\pi^2} + d_1^2\right)}
    \label{eq:barske_tip_diameter}
\end{equation}

Barske then gives these guidelines for other dimensions of the pump:

\begin{align}
    s_{ax} & \leq \min{\left(0.75, \frac{d_2}{100}\right)} \\
    b_1 & \geq \frac{1}{4} d_1 \\ %note to self, b is blade height
    b_2 & \geq \max{\left(b_1 \frac{d_1}{d_2}, 3s_{ax}\right)} \\
\end{align}

For the annular casing sizing, Barske doesn't have very specific numbers, so Gulich's guideline where B/H = 1.5 to 2.5 can be followed:

\begin{align}
    B & = 2 s_ax + b_2 \\
    H & = B / 2
\end{align}

\todo{look at efficiency that dynamicpumps uses, and cavitation maybe as well, iirc they have first principles}

To size the diffuser, Barske suggests sizing the diffuser throat to be around 0.8 to 0.9 of $u_2$, with the area of the diffuser outlet to be 3-4 times the area of the diffuser throat.

\subsection{Barske Pump efficiency}

Barske assumes that the disc friction losses can be approximated as a circular disc in a rotating fluid. Converting the equation into metric units gives:

\begin{equation}
    P_{RR} [\mathrm{W}] = 1956\times\rho\nu^{0.2}\left(\frac{n}{1000}\right)^{2.8}\left(d_2^{4.6}+4.6d_1^{3.6}b_1\right)
\end{equation}

As the hydraulic efficiency is defined as the useful power out of the pump over the total power into the pump, efficiency can be defined as:

\begin{align}
\eta_h & = \frac{\text{Useful power}}{\text{Input power}} = \frac{P}{P_{th}+P_{RR}}
 & = \frac{Qp}{Qp_{th}+P_{RR}}
 & = \frac{p}{p_{th}+P_{RR}/Q}
\end{align}

If total shaft efficiency is desired, the mechanical losses can be included:

\begin{equation}
    \eta=\frac{p}{p_{th}+\frac{P_F+P_{Mech}}{Q}}
\end{equation}

Lock provides a refined model which takes into account prerotation, nozzle losses and correction coefficients taking into account theoretical two dimensional flow solutions by A. Busemann \cite{buseman}. \todo{think if his 5 or so equations should be included here ot not, he also says that disc friction is treated as mechanical loss so doesn't effect head measurements and only effect efficiency calcluations, i will check if my code does this later on.}


By writing the total theoretical pressure gain in head:

\begin{equation}
    H_T = \frac{U_2^2 - U_1^2}{2g} + \frac{U_2^2}{2g}
\end{equation}

Lock uses busemann's corrections and derives a refined total head, taking into account prerotation:

\begin{equation}
    H_{TH} = h_0\frac{U_2^2}{g} - C_H\,K_1\!\left(1-\frac{Q}{Q_{op}}\right)\!\frac{D_1}{D_2}\,\frac{U_2^2}{g}
    \label{eq:lock_total_head}
\end{equation}

\begin{equation}
H_{SP} = H_{TH} - \underbrace{\zeta_{LD}\left[1-\left(\tfrac{D_3}{D_4}\right)^2\right]^2\frac{V_3^2}{2g}}_{\text{diffuser-nozzle loss}} - \frac{V_4^2}{2g} + \frac{V_0^2}{2g}
\label{eq:lock_static_head}
\end{equation}

Where $\zeta_{LD}=0.194$ is the diffuser loss coefficient from Hydraulic Institute's Pipe Friction Manual \cite{hydraulicpipemanual}.

By expressing $V^2/2g$ by $(8/g\pi^2)Q^2/D^4$ with $V=4Q/\pi D^2$, the head can be expressed as a function of flow rate:

\begin{equation}
    H_{SP} = h_0\frac{U_2^2}{g} - C_H\,K_1\!\left(1-\frac{Q}{Q_{op}}\right)\!\frac{D_1}{D_2}\,\frac{U_2^2}{g} - \frac{8Q^2}{g\pi^2}\left(\frac{\zeta_{LD}}{D_3^4} + \frac{1}{D_4^4} - \frac{1}{D_0^4}\right)
    \label{eq:lock_static_head_flow}
\end{equation}

To find the optimum flow rate used to shape the pump curve, the derivative of the head with respect to flow rate can be taken and set to zero:

\begin{equation}
    Q_{op} = \sqrt{\frac{h_0\frac{U_2^2}{g}}{\frac{24}{g\pi^2}\left(\frac{\zeta_{LD}}{D_3^4}\left(1-\left(\tfrac{D_3}{D_4}\right)^2\right) + \frac{1}{D_4^4} - \frac{1}{D_0^4}\right)}}
    \label{eq:lock_optimum_flow}
\end{equation}

Lock and Barske suggest that the flow cutoff point is limited by the diffuser throat. To find the static pressure at the diffuser throat, reversing the transfer between dynamic head into static head provided by the previous equations \todo{link to the previous equations}:

\begin{equation}
    H_3 = H_{SP} + \frac{8Q^2}{g\pi^2}\left(\frac{\zeta_{LD}}{D_3^4}\left(1-\left(\frac{D_3}{D_4}\right)^2\right) + \frac{1}{D_4^4} - \frac{1}{D_3^4}\right)
    \label{eq:lock_diffuser_throat}
\end{equation}

Then solving for the $Q$ at which $H_3$ is equal to vapor pressure head $H_{vap}=\frac{1}{\rho g}\left(p_{vap}-p_{inlet}\right)$ gives the flow cutoff point.

\todo{add cavitation of diffuser throat and shutoff flow here?}


\subsection{Turbine Literature and Theory}

\subsection{Impulse Turbine}

\todo{i dont know why reaction is bad here, oh and of course, need a diagram to show the velocity triangles.}

Turbines work with the same Euler turbine equation. In the supersonic impulse turbine design, the blades are only occupy a small outer portion of the blisk, so the turbine equation can be simplified:

\begin{equation}
    \Delta h_{\textrm{use}}=c_{3u}u_3-c_{4u}u_4=u(c_{3u}-c_{4u})
    \label{eq:turbine_specific_work}
\end{equation}

From \cref{eq:turbine_specific_work}, it is understood that a higher flow velocity difference $c_{3u}-c_{4u}$ per unit mass will lead to a higher $\Delta h_0$ according to equation \cref{eq:euler_specific_work}. This means that high absolute velocities, and for the maximum amount of the flow angled in the circumferential direction is optimal. For this reason, gas generator cycle turbines tend to take advantage of the high pressure ratios and thus the supersonic velocities. Then for the angles, the highest turning angle achievable is optimal. The amount of flow turning is usually limited by the achievable geometry.

\todo{use maybe sp125 to derive}

% todo: derivation

Leading to the ideal efficiency of an impulse turbine to be:

\begin{equation}
    \eta_{\textrm{ideal}} = 4\phi(\cos\beta_3-\phi)
    \label{eq:impulse_ideal_efficiency}
\end{equation}

Where $\phi$ is the flow coefficient defined as $\phi=c_3/u$ and $\beta_3$ is the blade angle at the inlet. The maximum ideal efficiency is achieved at $\phi=\cos\beta_3/2$. By taking the derivative of \cref{eq:impulse_ideal_efficiency}, the ideal $U/C_0$ is $0.5\cos\beta_3$. For a small turbine with a diameter of \todo{possibly add an example to show U/C0 for such turbines.}


% In my notes, this can be found "SP-8110.md" file and "Supersonic Turbine Blade Pitch and Height.md"
NASA special publication 8110 shows that for low $U/C_0$ turbines, impulse turbines are more efficient, even if reaction turbines can have a higher overall efficiency, just not at the low $U/C_0$ regime. In addition, the same conclusion is drawn by comparing the influence of tip clearance on the two types, where impulse turbines are more robust to tip clearance. This is especially important for small scale turbines where manufacturing tolerances are more difficult to achieve and contribute to more significant losses \cite{NASASP8110}. \todo{need to explain what reaction and impulse turbines are, maybe with a diagram?}



Another problem with designing small turbines is with the low blade height required. By calculating the mass flux through the turbine, the blade height can be designed:

\begin{equation}
    H_3 = \frac{\dot{m}}{\rho_3 c_{3m} d_{\textrm{mean}}}
    \label{eq:turbine_blade_height_full_admission}
\end{equation}

Where $\rho_3$ is the density of the turbine gas at the turbine inlet. Using the example turbine above, the required blade height would be \todo{addvalue}. By introducing partial admission, the mass flow through the turbine can be restricted to a smaller arc around the annulus at a slight cost to efficiency. The partial admission ratio $\zeta$ can be included into the blade height formula as so:

\begin{equation}
    H_3 = \frac{\dot{m}}{\zeta \rho_3 c_{3m} d_{\textrm{mean}}}
    \label{eq:turbine_blade_height}
\end{equation}

By reducing the partial admission to \todo{addvalue}, the blade height is increased to \todo{addvalue}. The design of the turbine blade height is very important due to multiple reasons. Firstly, the aspect ratio of the turbine blade determines the magnitude of of the wall boundary layer blockage. For this reason, a aspect ratio between \todo{addvalue} is desired.

Another way to increase the blade height is to decrease the meridonal velocity $c_{3m}$. This is achieved by increasing the angle that the flow has to turn as $c_{3m}=c_3\sin(\beta_3)$. 

For supersonic flow to develop within a turbine blade passage, a condition called "starting" must be satisfied. On the startup of a supersonic turbine, a normal shock starts at the inlet of the blade passage and as to pass through the whole passage before the flow into and out of the turbine passage can be supersonic. \todo{add derivation for this, later? also, loss models.}


\subsection{Loss models}
\label{subsec:weissloss}

The main loss model used in the design of the turbine is the one provided by Weiß \cite{weiss}. There are more general loss models provided by texts such as NASA SP-290, giving loss models for tip-clearance, disk friction, partial admission and incidence losses. However, those models were built on large full-admission, subsonic to transonic turbines. Weiss specifically derived his loss model for the specialized application of a supersonic impulse turbines at small scales (1-200kW). As the regime of the turbine considered in this thesis matches closer to Weiss's model compared to SP-290's model, Weiss's model is used. Weiss describes different models lead to stage efficiency differences of about 10\% points. Derived from rocket turbopumps literature and adjusted by in-house measurements, Weiss provides three efficiencies.


\begin{equation}
    k_N = \sqrt{1 - \left( 0.0029\,M_3^3 - 0.0502\,M_3^2 + 0.2241\,M_3 - 0.0877 \right)}
    \label{eq:weiss_nozzle_coeff}
\end{equation}

\cref{eq:weiss_nozzle_coeff} is the nozzle loss coefficient, relating the actual velocity to the theoretical isentropic velocity by $k_N=c_{No}/c_{No,is}$. 

\begin{equation}
    \begin{aligned}
        k_R ={}& 0.957
            - 3.62\times10^{-4}\,\Delta\beta
            - 2.58\times10^{-2}\,M_{w3} \\
            &+ 6.39\times10^{-6}\,\Delta\beta^2
            + 6.74\times10^{-2}\,M_{w3}^2 \\
            &- 7.53\times10^{-8}\,\Delta\beta^3
            - 4.30\times10^{-2}\,M_{w3}^3 \\
            &- 2.38\times10^{-4}\,\Delta\beta\, M_{w3}
            + 1.45\times10^{-6}\,\Delta\beta^2 M_{w3}
            + 4.25\times10^{-5}\,\Delta\beta\, M_{w3}^2
    \end{aligned}
    \label{eq:weiss_rotor_coeff}
\end{equation}

\cref{eq:weiss_rotor_coeff} is the loss of velocity along the blade passage through a theoretically impulse blade, relating the actual velocity to the theoretical velocity by $k_R=w_{exit}/w_{inlet}$.

\begin{equation}
    P_v = \frac{1.85}{2}\,(1 - \varepsilon)\,\rho_3 \left( \frac{n}{60} \right)^{3} d_m^4 \,(4.5\,b)
    \label{eq:weiss_ventilation}
\end{equation}

Finally, \cref{eq:weiss_ventilation} captures the losses due to partial admission. Weiss calls this loss ventillation losses, but it can also be called windage losses. Ventilation losses refers to the the non-admitted arc being dragged through stagnant air.

\todo{FUCKS SAKE LMAO, look at comment}
% Weiss says this: There are somewell-known turbine loss models in literaturee.g. by Refs. [39e41] etc. All of them are mainly developed for and based on the data from subsonic MW-class multistage steam and gas turbines working with full admission. 
% 39 - Traupel W Thermische, 40 - Craig HRM, 41 - Kacker SC


\subsection{Method of Characteristics Blade Profile Design}
\label{subsec:method_of_characteristics}

The design so far is based on 1D mean line theory. The design of the actual blade profile which will turn the supersonic flow in the desired direction is often accomplished using an application of the Method of Characteristics to transition a uniform flow field into vortex flow. \todo{add some history here, iirc sp290 has some stuff on why vortex flow, but here needs some stuff about how vortex flow was proven to be able to turn flows, and method of characteristics is used to turn the flow into such vortex flow, so obviously some explanation of vortex flow is required}. 

\subsubsection{Supersonic Turbine Blade Terminology}

Geometry is separated into the turbine blade, and the passage between the flow. Two surfaces are separately designed, and it is common to refer to the top surface as the convex, and the bottom surface as the concave one. This is in relation to the solid blades themselves as opposed to an alternate designation by the top and bottom of the blade passage.

Vortex flow is a special type of supersonic flow in which flow follows a theoretical perfect arc around a central point. This is a very special type of flow field that is achieved by conditions where expansion and compression waves cancel each other to allow shockless rotation of the flow field. The vortex flow is bounded by the two circular arcs of the surfaces of which a desired surface Mach number is achieved. MOC generates the profile for desired mach numbers. Later analysis will need to be required to choose these desired Mach numbers.

A special condition of vortex flow satisfies the equation

\begin{equation}
    \frac{V}{r} = \textrm{constant}
    \label{eq:vortex_condition}
\end{equation}

Boxer and Goldman nondimensionalizes this equation using the critical velocity $M^{*}=V/a^*$ where $a^*=\sqrt{\gamma R T^*}$ is the speed of sound at the sonic point. As $a^*$ only depends on stagnation temperature, it is constant along any streamline of fixed stagnation enthalpy. Then, a nondimensional radius $r^*$ is introduced as the sonic velocity streamline in the vortex field, giving the following condition:

\begin{equation}
    \frac{M^*}{r^*} = 1
    \label{eq:vortex_condition_nondim}
\end{equation}

The Prantdl-Meyer angle can be defined from the critical mach number $M^*$ as:

\begin{equation}
\nu(M) = \frac{\pi}{4}\left(\sqrt{\frac{\gamma+1}{\gamma-1}} - 1\right) + \frac{1}{2}\left(\sqrt{\frac{\gamma+1}{\gamma-1}} \arcsin\left[(\gamma-1)M^{*2} - \gamma\right] + \arcsin\left(\frac{\gamma+1}{M^{*2}} - \gamma\right)\right)
\label{eq:prandtl_meyer_angle}
\end{equation}

And the converting between Mach number and Critical Mach number is given by:

\begin{equation}
M^* = \sqrt{\frac{(\gamma+1)M^2/2}{1 + (\gamma-1)M^2/2}}
\label{eq:mach_to_critical_mach}
\end{equation}

Equations \cref{eq:vortex_condition_nondim} and \cref{eq:mach_to_critical_mach} can be used to derive the nondimensional radius of the upper and lower surfaces of the blade. The end points of this circular arc is exactly the angle by which the inlet flow has to be turned to achieve the surface mach number. 

\begin{equation}
\begin{aligned}
    \alpha_{l-i} & = \beta_1 - (\nu_i - \nu_l) \\
    \alpha_{u-i} & = \beta_1 + (\nu_i - \nu_u)
\end{aligned}
\label{eq:inlet_flow_turning_angle}
\end{equation}

The equations for the outlet are the same as in \cref{eq:inlet_flow_turning_angle}.

\subsubsection{Transition arcs}

To transition from the uniform inlet flow to the vortex flow, method of characteristics is used to generate the blade profile to achieve this in a shockless manner. 

\todo{explain new axis}

To achieve this, the flow angles from $\beta_i$ to $\alpha_i$ is discretized into $k$ interval of $\phi_0$ to $\phi_k$. It is shown by \todo{find the ref from goldman 1968 ref 2} that the flow angle and dimensionless radius $R^*$ is related along its characteristic line with:

\begin{equation}
    \phi = \pm \frac{1}{2}f(R^*) + \textrm{constant}
    \label{eq:characteristic_line_equation}
\end{equation}

The positive gives the expansion line and the negative gives the compression line. Where $f(R^*)$ is:

\begin{equation}
    f(R^*) = \sqrt{\frac{\gamma + 1}{\gamma - 1}}\arcsin\left(\frac{\gamma-1}{R^{*2}} - \gamma\right) + \arcsin\left(\left(\gamma+1\right)R^{*2} - \gamma\right)
    \label{eq:f_R_star}
\end{equation}

$f(R^*)$ for a given flow angle can be calculated with:

\begin{equation}
    f(R^*_k) = 2\nu_i - \frac{\pi}{2}\left(\sqrt{\frac{\gamma+1}{\gamma-1}} - 1\right) - 2k\Delta\nu
    \label{eq:f_R_star_k}
\end{equation}

The result of \cref{eq:f_R_star_k} can be used to numerically solve for $R^*_k$ with \cref{eq:f_R_star}, and points corresponding to each flow angle along the major expansion characteristic line can be calculated with:
\begin{equation}
    \begin{aligned}
    x^*_{k} & = -R^*_k \sin(\phi_k) \\
    y^*_{k} & = R^*_k \cos(\phi_k)
    \end{aligned}
\end{equation}

The wall segments are generated by starting from the end point of the circular arc, then finding the intersection of the previous wall segment and the mach line from this point. The slope of the mach line with respect to to the $x^*$ axis is given by:

\begin{equation}
    m_k = \tan\left(\frac{\phi_k+\phi_{k-1}}{2}+\frac{\mu_k+\mu_{k-1}}{2}\right)
\end{equation}

Where
\begin{equation}
    \mu_k = -\arcsin\left(\frac{1}{M_k}\right) = -\arcsin\left(\sqrt{\left(\frac{\gamma+1}{2}\right)R^{*2}_k - \frac{\gamma-1}{2}}\right)
\end{equation}

By iterating through the flow angles for both the upper and lower surfaces for the inlet and outlet, the full theoretical MOC blade profile can be generated.

\todo{add "i" to the variables to define inlet}

\subsection{Choosing mach numbers}

There are two main factors to consider when choosing the inlet, outlet and surface mach numbers. One is to prevent flow separation, and the other is to ensure that supersonic flow within the blade passage can "start".

\subsubsection{Flow separation}

Separation occurs in high adverse pressure gradients. In supersonic flow, this happens whenever the flow is decelerating along the surface.  There are two main regions on the blade surfaces which separation can occur. At the upper surface, the inlet transition arc accelerates the flow from the inlet velocity until the vortex flow, then decelerates the flow until the outlet. At the lower surface, the inlet transition arc decelerates the flow from the inlet velocity until the vortex flow, then accelerates the flow until the outlet. 

\todo{add a figure, question: literature figure or my figure? probably both is best}

Goldman believes that separation on the lower surface is probably not as important as the flow would tend to reattach shortly as the pressure gradient becomes favourable again \cite{goldman1970}. Goldman uses the value of the incompressible form factor $H_i$ along the blade surface to predict separation, where  Schlichting found \todo{find goldman ref 12} that separation usually occurs in the range of $H_i$ from 1.8 to 2.4.

To find $H_i$ along both surfaces, Sasman and Cresci \cite{sasmancresci} provides a method to calculate the compressible turbulent integral boundary layer with two ODEs. 

Once a MOC blade profile is generated and the mach number at every point along the blade is known, the momentum-integral equation and shape-factor equation can be solved as an initial value problem starting from the leading edge goverened by the following equations:

\begin{equation}
    \frac{df}{d(s/c)} = 1.268\left[-\frac{f}{M_e}\frac{dM_e}{d(s/c)}\left(1 + g_w H_i\right) + A\,c\right]
\end{equation}

\begin{equation}
    A\,c = 0.123\,e^{-1.561 H_i}\, M_e\, \mathrm{Re}_0\, \frac{T_e}{T_0}\left(\frac{\bar\mu}{\mu_0}\right)^{0.268}
\end{equation}

\begin{equation}
\frac{dH_i}{d(s/c)} = \underbrace{-\frac{1}{2M_e}\frac{dM_e}{d(s/c)}\,H_i(H_i+1)^2(H_i-1)\left[1 + (g_w-1)\frac{H_i^2+4H_i-1}{(H_i+1)(H_i+3)}\right]}_{\text{pressure-gradient term: deceleration drives } H_i \text{ up}} \\ + \underbrace{\frac{H_i^2-1}{f}\,A\,c\left[H_i - \frac{0.011\,(H_i+1)(H_i-1)^2}{H_i^2\,(C_f/2)}\frac{T_e}{\bar T}\right]}_{\text{shear term: wall friction relaxes } H_i \text{ back toward equilibrium}}
\end{equation}

With initial conditions of a flat-plate $\frac{\theta_0}{c} = 0.036\,\frac{s_0}{c}\,\mathrm{Re}_{s_0}^{-1/5}$. \todo{cite this}
% i found delta = 0.37x/(Re to 1/5) from anderson 19.1 page 1052, prandtl u/U = (y/delta) to 1/7 then shoving into the equation of momentum thickness gives 7/72 delta, which is 0.0359722222, done.

The initial condition for the shape factor is chosen so that it matches closely to the rapidly converging value $ H_{i,0} = 1.67$. \todo{make this more rigorous}

It is important to note that Goldman provides equations in his NASA TN D-4421 1968 paper that has been shown to not be conservative enough, with tests by Erikson saying that his criteria is a little too optimistic. The method described here is from Goldman's NASA TM X-2095 paper.

\todo{reorganize here}

\todo{should I remove these?}
%  it took me a while to understand, but I think I finally understand it. Sasman and cresci developed a way to transform (stretch) the surface arc coordinates from compressible to incompressible coordinates, somehow deriving their eq 18, 19, 20 which can be solved as an IVP with initial conditions in 

\todo{addd graphs to show validation against goldman, I already have this.}

\subsubsection{Starting}

On startup before the turbine passage runs supersonically, a detached bow shock stands ahead of the leading edges. The flow behind this normal shock is subsonic and has lost a significant amount of total pressure. For the turbine to operate nominally with supersonic flow, the passsage must be able to "swallow" this shock, allowing the shock to move through the whole passage.

To do so, the do so, the post shock flow with a reduced total pressure must have enough of a flow area pass the required mass flow similar to the choked flow condition. A similar condition is the unstart of a supersonic wind tunnel or a supersonic internal intake. However, Boxer noticed in 1952 that due to the large velocity gradient between the channel walls due to tue turning, the flow cannot be assumed one directional as in the previous examples. Goldman in 1968 provides the condition accounting for the "reduction of mass flow due to curvilinear flow" to calculate the maximum inlet design mach number to ensure starting as:

\begin{equation}
\frac{p_{02}}{p_{01}}(M_i) \;\geq\; \frac{Q}{1-C}
\end{equation}

Where $p_{02}/p_{01}$ is the total pressure loss across the normal shock, $M_i$ is the inlet design mach number, $Q$ is named the vortex parameter \cite{sudhof} and $C$ is the reduction in mass flow due to two-dimensional flow. Solving this equation for $M_i$ gives the maximum inlet design mach number to ensure starting. \todo{perhaps put in the starting graph, showing starting limit as a function of inlet and surface mach numbers?}

Historically, starting was a major problem for supersonic turbines. Before the theory was established, Colclough had to build a special jig in which every blade could be pivoted to vary the contraction until the blade started.

\subsubsection{Mach number selection}

The separation criterion limits the acceleration of the flow along the upper surface, and the deceleration of the flow along the lower surface. Then, the starting criterion provides an additional limit on the inlet design mach number based on these surface mach numbers. 

However, Goldman notes that for optimum efficiency, the maximum spread of surface number leads to maximum efficiencies. It comes from the obvious fact that a loss of velocity across the impulse blade leads to a loss of kinetic energy, so he posits the efficiency of an impulse blade to be:

\begin{equation}
    \eta = \left(\frac{w_{exit}}{w_{exit,ideal}}\right)^2
\end{equation}



This is nearly the same as Weiss' \cref{eq:weiss_rotor_coeff} loss. Weiss derived his empirical model only of 1D parameters such as inlet mach number and the angled the flow is turned. 

What this efficiency tries to capture is the boundary layer friction losses, shocks, and exit wake mixing losses. The effect of spreading the surface mach numbers becomes clear by looking at the vortex equation \cref{eq:vortex_condition_nondim}. The larger the mach number spread, the larger the difference in $R^*$ between the upper and lower surface, the larger physical passage width. The wider the channel compared to the blade chord, the lower fraction of the boundary layer thickness to passage width, the lower the wetted surface area per unit flow, and thus the lower the losses.

This is the challenge of supersonic turbine blade design, maximum efficiency is achieved by riding along the edge of separation and starting.

\todo{two MOC channels with small spread and large spread, eff too?}

\subsubsection{Finite leading edge}

In \cref{subsec:method_of_characteristics}, the upper and lower surface are mathematically designed to produce no shocks in transitioning the unifrom flow from the nozzle exit into the vortex flow. For no shocks to form, both surfaces must be perfectly parallel to the starting flow. This is problematic because this implies the point where the two surfaces intersect at the leading edge, the angle between the upper and lower surfaces is zero. This creates an impossible to manufacture zero thickness leading edge. Boxer suggests a few methods to solve this issue. He separates design into turbines with a subsonic or supersonic axial velocity component. 

For subsonic axial velocity turbines, either lower or upper surface can be designed with an incidence angle to the inlet flow. Angling the upper convex surface and keeping the lower concave surface parallel to the inlet flow, the compressive oblique shock will be generated along the region between the inlet and the nozzle exit, affecting the flow in the next blade passage. Kantrowitz shows that in steady state a compression shock of equal strength will be created to obtain the steady state condition \cite{kantrowitz}. To solve this, a small angle wedge can be placed shortly after the leading edge to generate a small expansion wave to cancel out the shock. 

Another method is angling the lower concave surface so that the shock stays within the blade passasge and does not affect the upstream flow. The oblique comrpession shock can then be designed to intersect the upper surface at an angle to fully cancel out the shock.

For supersonic axial velocity turbines, a wedge can be designed on the convex surface with a low enough angle from the design inlet flow so that the shock stays within the blade passage and can be cancelled out by a kink on the concave surface.

\todo{figures to explain this}


\subsubsection{Nozzle Design}

Historically, the optimal nozzle design for a ideal inviscid gas was designed with the Method of Characteristics as well. Rao provides a method to achieve this \cite{raonozzledesign}. NASA special pubplication 8120 explains that a single canted parabola closely approximates the ideal nozzle contour and should be used for high performance nozzles. In this paper, the nozzle was a simple conical nozzle with a 7 degree half angle. It can be seen from \todo{add figure 7 from nasa sp8120} that the performance of a conical nozzle for low area ratios lead to a efficiency loss below 2\%.

\section{Turbomachinery Design}

The design point for this thesis was chosen from the capabilities of the test rig. The design parameters for the pump and turbine is shown in \todo{addcreftable} and this section explains how each number was chosen.

The RPM limit was set by the maximum speed of the Magtrol TM306 inline torquemeter and the AHB-3 hysteris brake at 20,000 RPM. The pressure rise and flow rate was chosen to study the low specific regime. The supply pump was decided to be Hozelock's JET 3000 K7 Jet pump as it was available in the department and had a published performance curve. At \valPumpMdotDesign, the supply pump had a head of around \valSupplyPumpHeadDesign. Then, the head rise of the pump was chosen to be \valPumpPressureDesign for a specific speed of \todo{addspecificspeed}.



% \todo{does lock's model change barske's design methodology for d2? I can write the barske model for now and say future work can look at rearranging lock maybe?}

\section{Turbine Sizing and Mechanical Design}


\subsection{Turbine Design Point}

Turbine sizing has to satisfy many requirements simultaneously, many design variables. After the pump was sized for \valPumpRpmDesign at a power of \valTurbPowerDesign, there are many remaining free variables.

The main parameters being turbine diameter $D_mean$, mass flow through the turbine $\dot{m}$, blade inlet angle $\beta_3$.

\section{Experimental Setup}

The experimental setup was designed for three specific configurations: Electric motor driven pump testing, Isolated turbine coupled to a hysteris brake, and integrated turbine and pump testing.

\subsection{Test Configurations}

\begin{figure}[H]
    \centering
    \begin{subfigure}[b]{0.3\textwidth}
        \centering
        
\includegraphics[width=\textwidth]{placeholder.png}
        \caption{TODO}
        \label{fig:electric-pump-config}
    \end{subfigure}
    \hfill % Adds flexible space between figures
    \begin{subfigure}[b]{0.3\textwidth}
        \centering
        
\includegraphics[width=\textwidth]{placeholder.png}
        \caption{TODO}
        \label{fig:turbine-brake-config}
    \end{subfigure}
    \hfill % Adds flexible space between figures
    \begin{subfigure}[b]{0.3\textwidth}
        \centering
        
\includegraphics[width=\textwidth]{placeholder.png}
        \caption{TODO}
        \label{fig:turbopump-config}
    \end{subfigure}
    \caption{TODO}
    \label{fig:test-configs}
\end{figure}

In all three configurations, there is a driving component, the torquemter and a load. \cref{fig:electric-pump-config} shows the electric motor driven pump testing configuration. The electric motor test is used to gather most of the data because the motor is more controllable, repeatable and safer. In this configuration. The second configuration in \cref{fig:turbine-brake-config} connects the turbine to a hysteris brake to test the turbine in isolation. Due to time constraints, this configuration was not tested. The final configuration in \cref{fig:turbopump-config} connects the turbine and pump together to test the performance of the integrated system.

%torqumeter accuracy chain: dsp6001 0.05% VDC accuracy, input from tm306 at 0.02% accuracy, from <0.1% accuracy in tm306 itself. 
\subsection{Mechanical Layout and Shaft Train}

At the design 20,000 RPM for the rig, shaft alignment, rigidity and maintainability were the main drivers. At 20,000 RPM, available shaft couplers had barely any tolerance for misalignment. The Magtrol MIC-5-2470 couplers had a maximum 0.8mm axial, 0.36mm radial and$0.7^\circ$ angular misalignment tolerance at 35,000 RPM. Shaft couplers had to be used to connect rotary components to the transducers, a requirement that doesn't have to be satisfied by turbopump units as the shaft can be machined as a single piece.

The central component of each configuration in \cref{fig:test-configs} is the Magtrol's TM306 inline torquemeter as it allows the stated <0.1\% torque measurement accuracy. As the torquemeter was designed to be mounted onto a base plate, a base plate design was chosen instead of an integrated assembly like a normal turbopump.

The base plate is a \valRigBaseplateThk plywood with a laminate on both sides to be semi-waterproof, to have a flat surface for mounting, rigid but also cheap. M6 and M8 Countersunk holes were drilled into the base plate with a CNC router to allow for precise mounting of the components. 

\subsubsection{Shaft alignment}

The reason the shafts had to be aligned was because flexible shaft couplers were used. These tolerate some misalignment at high speeds, an essential requirement for a test rig where two or more shaft have to be mechanically connected.

The alignment of each component to the torquemeter had many challenges. Mounting magnetic dial indicators onto the shafts was not possible due to three reasons. The shafts were made of austenitic stainless steels so they were non-magnetic, the shafts were short and had a low diameter (10-15mm diameter, with maximum stub length of 20mm). A short stub length is highly desirable in rotordyanmics to increase rigidty and increase the first critical speed.

The shaft couplers were also single piece, so the shaft couplers which do not have a edge true to the shaft axis. This meant that the shaft couplers could not be installed after the shafts were aligned, as the couplers would have to be installed before the shafts were aligned. The space the shaft couplers took up meant it was even more difficult to use the shaft stubs for alignment.

It was decided to use the flat faces of the torquemeter base as the reference for alignment. By designing the pump and turbine mounts to have a similar rectangular base and assuming that each shaft was true to the rectangular base, alignment was finally possible. Angular alignment was achieved by 3D printing a spacer with two parallel edges for the pump or turbine to sit flat against onto the torquemeter. Then, radial misalignment was achieved by designing the face parallel to the shaft to be the parallel to the torquemeter face. A straight edge could then be used as a guide to align the shaft while tightening the mounting bolts. This procedure is shown in \cref{fig:shaft-alignment}.

\begin{figure}[H]
    \centering
    
\includegraphics[width=0.5\textwidth]{placeholder.png}
    \caption{TODO}
    \label{fig:shaft-alignment}
\end{figure}

As the TM306 torquemeter's shaft with the base plate was machined to be preciscely 80.00 mm from the base plate, this same height was built into the pump and turbine mounts to achieve radial aligment in the vertical direction.


\subsection{Water Circuit}

The water circuit is as follows: A water tank is pumped by the supply pump into an inlet throttle valve then into the experimental pump inlet. The pump outlet is then throttled by an outlet throttle valve before being returned to the water tank. Pressure sensors are placed at the inlet and outlet of the pump to measure the head rise across the pump. Flow rate is measured by a flow meter placed in the inlet line before the inlet throttle valve instead of the return line due to space constraints placed by the blast shield.

\subsubsection{Inlet and Outlet Throttle Valves}

The inlet throttle valve allows a reduction in pump inlet pressure to characterise the cavitation behaviour of the pump. The outlet throttle valve allows for a reduction in flow rate to characterise the head and cavitation behaviour of the pump at different flow rates.

The selection of the size of the inlet and outlet valves are crucial to the quality of the data. Ball valves were chosen as there were spares available, and they could be throttled in a controllable and repeatable manner by the angle of the valve handle. Custom 3D printed housings were designed to attach a servo motor to the valve handle in direct drive to allow for precise, remote and repeatable control of each valves. 

The main requirement for both valves were so that the range of required flow rates and pressure drop through each valve corresponded to a large range in valve opening angles. For the inlet valve, fully open must have a large margin to the cavitation limit of the pump, ideally, cavitation would start at around 50\% valve opening and could be further reduced until the valve is closed. For the outlet valve, fully open must be at least 1.5 times the cut-off flow rate of the pump.


%cite:
% https://www.mydatabook.org/fluid-mechanics/flow-coefficient-opening-and-closure-curves-of-full-bore-ball-valves/
% https://www.engineeringtoolbox.com/ball-valves-flow-coefficients-d_223.html
% https://tameson.co.uk/pages/ball-valve-size-calculator

This requirement is further driven by the $1-3^\circ$ resolution of the servo motors available in the department. Data for Kv against opening angle were extrapolated from online data \todo{add citation}  from 1/8" in to 3/4" in were studied. However, as the data was extrapolated from data of ball valves with larger diameters, they had to be experimentally validated. The flow through a ball valve is also highly dependent on inlet and oulet conditions such as contraction or expansion of pipe diameter at the ports \cite{swagelokMS-06-84}. The size of my pipes were limited to what I could find in spare drawers, adding to the uncertainty.

\begin{figure}[H]
    \centering
    \begin{subfigure}[b]{0.4\textwidth}
        \centering
        
\includegraphics[width=\textwidth]{placeholder.png}
        \caption{TODO}
        \label{fig:servo-valves}
    \end{subfigure}
    \hfill % Adds flexible space between figures
    \begin{subfigure}[b]{0.4\textwidth}
        \centering
        
\includegraphics[width=\textwidth]{placeholder.png}
        \caption{TODO}
        \label{fig:valve-kv-test}
    \end{subfigure}
    \caption{TODO}
    % \label{fig:TODO}
\end{figure}

\cref{fig:valve-kv-test} shows the experimentally measured Kv against opening angle values were lower than expected. However, this worked in my favour as it actually increased the range of useful valve opening angles.

The location chosen for the experiment required a safe test area and a control room separated to allow for the failure of the test rig without risk to the operator. For this reason, the experiment was conducted within the Imperial College Hypersonics Lab. The test room contains equipment that had to be protected, so a blast shield was designed and set up around the turbomachinery stand to protect the equipment from the possibility of mechanical failure and fragmentation.

\subsection{Air supply}

Originally, a 100 bar compressed air supply was designed and parts readied to be used to drive the turbine. Unfortunately due to time constraints, the shop air supply within the hypersonics lab was used instead. Details of the 100 bar compressed air supply can be found in \todo{add appendix reference}. The shop air which was limited to 7 barg was connected to partially assembled parts of original feed system. The shop air was connected to a GRT20s-280N100-SSKN-L 0-280 to 0-100 bar regulator with a Kv of 1.5. This regulator fed into the the main 1" valve, teed into a 1/2" vent valve and finally into the turbine inlet.

Both valves were Swagelok pneumatically actuated spring return ball valves. The main valve was the  SS-65TS16-35C with a Cv of 40 and the vent valve was SS-45S8-33C with a Cv of 12. These valves were driven by 2 5/2 solenoid valves connected to the shop air supply.

Due to droop in the shop air regulator, maximum turbine inlet pressure at max flow was measured to be 3.1 bara. There was no budget for an air flowmeter.

\subsection{Instrumentation and Data Acquisition}

For cost and time reasons, a custom data acquisition and control system was designed to allow for remote operation of the test rig. 

\subsubsection{Main Architecture}

For safety reasons, the electronics were enclosed in a plastic box placed in the test room, controlled by a LAN cable fed under the wall between the test room and the control room into my laptop.

The main main computer inside this box was a Raspberry Pi 5 running a backend written in Python. The Raspberry Pi communicated with 3 ESP32S3 microcontrollers over a UART serial connection at 921600 baud. One ESP32S3 was used as the actuator controller, one was used as a ADC and flowmeter DAQ, and one was used as a thermocouple DAQ. Breadboards would require too many wires to maintain and unstable, while there was not enough time to make custom PCBs. Soldered perfboard was used to connect the components. The Raspberry Pi communicates through the LAN cable through UDP to a laptop in the control room, allowing for live data readout and control of the test rig with a GUI written in Python and Qt.

\subsubsection{ADC Measurements}

Pressure, torque, RPM, voltage and current were measured with a MCP3208 8 channel 12-bit ADC connected to the ESP32S3. The MCP3208 was powered by a LT3042 LDO at 5.11V to reduce noise. This 5.11V was also used to power the pressure sensors. As the pressure transducers were ratiometric, both the ADC and the pressure transducers track bacck to the same voltage reference, reducing drift and noise.

The custom backend could only send UART signals back and forth to the ESPs at a 50 Hz loop. The MCP3208 was capable of sampling up to 100 ksps, so the ESP32 was programmed to collect 32 samples and average them per channel before sending the averaged data back to the Raspberry Pi. High frequency data couldn't be collected due to the UART connection being the bottleneck, but this averaging allowed for measurement of voltage at a resolution higher than 12-bits.

Originally, a separate computer running Magtrol's M-Test software was planned to be used to read the torque and RPM from the DSP6001. However, by connecting the DSP6001's analog output for torque and RPM to the MCP3208, the torque and RPM data could be integrated into the same DAQ system as the pressure sensors, allowing for time-exact data synchronisation and simpler data collection.

To measure the torque and RPM, the TM306 torquemeter was connected to the Magtrol DSP6001 dynamometer controller. The DSP6001 outputs a -10V to 10V analog signal for torque, and 0 to 10V for RPM. The DSP6001 was set to output 1V per 0.25 Nm and 1V per 5000 RPM. Because the MCP3208 could only measure positive voltages, an OPA2134 op-amp was used to shift the expected -5V to 5V torque signal to 0V to 5V. As the op-amp was not rail-to-rail, a buck converter was used to generate a -9V and 9V supply for the op-amp.

The motor's voltage and current were also measured by the ADC. The 24V supply to the motor ESC was measured with a voltage divider, and the current was measured with a WCS1600 DC hall effect current sensor.

\subsubsection{Other sensors}

K-Type thermocouples although are less accurate for such a small temperature range were chosen for their cost. The thermocouples were read by an Adafruit MAX31855 thermocouple amplifier breakout board with integrated voltage regulator connected to the ESP32S3. 

A E-WFS-GE turbine flowmeter with a range of 2 to 45 L/min with a cited accuracy of $\pm10\%$ was used. To calibrate the flowmeter, a bucket and stopwatch was used to compare the flowmeter to the integrated flow rate over a period of time. The flowmeter was found to be accurate to $\pm3\%$ \cite{flowmtr}. The turbine generates pulses based on its rotation rate. The ESP32 measures the interval between these pulses to calculate the flow rate. % https://www.networktechinc.com/pdf/e-wfs-g.pdf

\subsection{Actuator Control}

This ESP32 was responsible for controlling the inlet and outlet throttle valve angles, 

\section{Results and Discussion}

\todo{the lower head might be due to leading edge rouding as said by lock, maybe not due to prerotation as the equations suggest.}

\todo{look at H-Q data and see if the cavitation boundary lock describes can be found there as my current cavitation curves look at a different thing.}

\subsection{Experimental Campaign}

The experimental campaign was conducted in three phases: commissioning, electric motor driven pump testing, then integrated turbine and pump testing. The standalone turbine testing was not conducted due to time constraints.

\subsection{Pump results}

There are two main curves that are used to evaluate pump performance: the head-flow curve and the efficiency curve. The head-flow curve shows the achievable head at different flow rates and will be discussed first.


% \begin{figure}[H]
%     \centering
%     \begin{subfigure}[b]{0.4\textwidth}
%         \centering
%         
\includegraphics[width=\textwidth]{placeholder.png}
%         \caption{TODO}
%         % \label{fig:TODO}
%     \end{subfigure}
%     \hfill % Adds flexible space between figures
%     \begin{subfigure}[b]{0.4\textwidth}
%         \centering
%         
\includegraphics[width=\textwidth]{placeholder.png}
%         \caption{TODO}
%         % \label{fig:TODO}
%     \end{subfigure}
%     \caption{TODO}
%     % \label{fig:TODO}
% \end{figure}

% \begin{figure}[H]
%     \centering
%     
\includegraphics[width=0.5\textwidth]{placeholder.png}
%     % \label{fig:TODO}
% \end{figure}

% \begin{figure}[H]
%     \centering
%     \begin{subfigure}[b]{0.4\textwidth}
%         \centering
%         
\includegraphics[width=\textwidth]{placeholder.png}
%         \caption{Dimensional head-flow curve}
%         \label{fig:dim-hq}
%     \end{subfigure}
%     \hfill % Adds flexible space between figures
%     \begin{subfigure}[b]{0.4\textwidth}
%         \centering
%         
\includegraphics[width=\textwidth]{placeholder.png}
%         \caption{Non-dimensional head-flow curve}
%         \label{fig:non-dim-hq}
%     \end{subfigure}
%     \caption{TODO}
%     % \label{fig:TODO}
% \end{figure}

\begin{figure}[H]
    \centering
    \includegraphics[width=1\textwidth]{theory_HQ_tuned.pdf}
    \label{fig:dim-hq}
    \caption{TODO}
\end{figure}


\todo{add explanation that RPM was not constant, so affinity laws were used to adjust the data to the average rpm}

% \cref{fig:dim-hq,fig:non-dim-hq}

\cref{fig:dim-hq} Shows the actual achieved head rise across flow rates at the different RPMs that were tested. Overlayed are solid and dashed lines. Both are theoretical head-flow curves generated from \cref{eq:lock_static_head_flow}. The dashed line is with exactly the parameters from Lock's paper, and the solid line is with the pre-rotation $K$ increased from \valLockKOneOriginal to \valLockKOneModified. The dashed vertical dashed line shows the point at a synchronized video camera detected cavitation within the diffuser through the transparent diffuser wall. It can be seen that the original parameters overpredict the head rise, and there are many mechanisms that could be causing this. 

The efficiency curve generated by measuring shaft power, subtracting mechanical losses can be shown in the following graph:

The right graph shows the predicted efficiency curve from pure theory, and the right graph shows the curve but with adjusted disc friction. To explain the disc friction adjustment, the pump data at zero flow where all the shaft power was only going into hydraulic losses:

\begin{figure}[H]
    \centering
    \begin{subfigure}[b]{0.49\textwidth}
        \centering
        \includegraphics[width=1\textwidth]{results_eta_flow.pdf}
        \label{fig:eff}
        \caption{TODO}
    \end{subfigure}
    % \hfill % Adds flexible space between figures
    \begin{subfigure}[b]{0.49\textwidth}
        \centering
        \includegraphics[width=\textwidth]{results_churning_power.pdf}
        \caption{TODO}
        \label{fig:churning}
    \end{subfigure}
    \caption{TODO}
    % \label{fig:TODO}
\end{figure}

\cref{fig:churning} shows that the measured disc friction losses are about \valPumpChurnMult times higher than barske's formula for disc friction. Adjusting the disc friction by this factor leads to a much better fit of the efficiency curve, as can be seen in \cref{fig:eff}.

Specifically for rocket engine applications, the cavitation behaviour is of high importance. The following graph shows the cavitation curve generated by plotting the head at which cavitation was observed at different flow rates:


\begin{figure}[H]
    \centering
    \includegraphics[width=0.7\textwidth]{results_npshr_vs_q.pdf}
    \caption{TODO}
    \label{fig:npshr}
\end{figure}

   
\cref{fig:npshr} shows the NPSH required to avoid cavitation at different flow rates. The dashed line is the theoretical curve generated by solving \cref{eq:lock_diffuser_throat} for the flow rate at which the head at the diffuser throat is equal to the vapor pressure head. Two interesting observations here.


\todo{add picture of no cavitation, cavitation of throat only, cavitation of inlet, and both}
\todo{3\% vs full cavitation?}